\section{Experimental design}
\label{sec:experiments}

The experiments ask whether rankings produced by CIT and CIF place predictive
features early enough for fixed downstream learners, and what the
conditional-inference controls cost. For each training fold, a method produces
a feature ranking, and fixed downstream learners evaluate the selected features
on the held-out fold.

\subsection{Benchmark setup}
\label{sec:experiments-protocol}

The benchmark uses 5 random seeds, with 5-fold cross-validation within each
seed. Classification fits logistic regression (LR), support vector machine
(SVM), and k-nearest neighbors (KNN) as downstream models and reports balanced
accuracy, defined as mean recall across classes. Regression fits ridge
regression, support vector regression (SVR), and KNN and reports $R^2$. The main
tables use the standard values $k \in \{5,10,25,50,100\}$.
Downstream learner settings are fixed. Logistic regression uses balanced class
weights and a 1,000-iteration limit; SVM uses an RBF kernel with unit
regularization, kernel scale set by the \texttt{scale} rule, and balanced class
weights; KNN uses 5 distance-weighted neighbors; ridge uses unit penalty
strength; and SVR uses an RBF kernel with unit regularization, $\epsilon=0.1$,
and the same kernel scale rule.

For each seed and fold, feature-standardization parameters are estimated on the
training fold and then applied to the held-out fold. The ranking method uses only
the standardized training fold and produces the ordered feature list described in
\cref{sec:rank-construction}. For each $k$, we take the top-$k$ features from
that ranking, refit standardization on the selected training columns, apply the
same transform to the selected held-out columns, fit the downstream model on the
selected training data, and evaluate it on the held-out fold.

The downstream evaluation also records performance using the full feature set,
$k=p$, for every dataset. For datasets with more than 100 features, the separate
high-$p$ analysis adds fixed values $k \in \{150,200,300,500,750,1{,}000\}$
and fractional values
$\lceil 0.25p\rceil,\lceil 0.5p\rceil,\lceil 0.75p\rceil$, keeping only values
no larger than $p$. The main tables average only the standard $k$ values.

\subsection{Benchmark comparisons}
\label{sec:experiments-analysis-plan}

The benchmark averages scores in three steps. First, for each
$(\text{dataset}, \text{downstream model}, k)$ combination, scores are averaged
over folds and seeds. Second, each dataset score averages those fold-and-seed
means over downstream models and over the standard $k$ values available for the
methods in that comparison. Third, the tables average dataset scores over
datasets.

Some method runs on \nolinkurl{gisette}, \nolinkurl{isolet}, and
\nolinkurl{orlraws10P} were excluded (\cref{app:methods}), so the all-method
classification panel has 21 of the 23 datasets, and each pairwise comparison
uses the datasets where both methods are present. Two smaller complete-case
panels, used for bootstrap intervals, omnibus rank tests, and LODO checks,
retain only datasets with every method, every downstream learner, and all five
standard values of $k$. A dataset with fewer than $k$ features does not enter
the column for $k$, so the classification counts by $k$ are 21, 20, 14, 14, and
13 and the regression counts are 8, 8, 7, 7, and 6. \Cref{tab:panels} names
the panels.

\begin{table}[!htbp]
  \centering
  \footnotesize
  \setlength{\tabcolsep}{4pt}
  \begin{tabular}{@{}lp{4.3cm}rrp{4.6cm}@{}}
    \toprule
    \tablehead{Panel} & \tablehead{Definition} & Classification & Regression & \tablehead{Used for} \\
    \midrule
    All-method & every method present & 21 & 8 & main rank table, LODO \\
    Pairwise & both compared methods present & 21 to 23 & 8 & direct CIF comparisons \\
    Complete-case & every method, learner, and $k$ present & 13 & 6 & bootstrap intervals, omnibus tests, LODO \\
    High-$p$ & real datasets with more than 100 features & 15 & 6 & high-$p$ analysis \\
    Synthetic & datasets with known informative features & 8 & 8 & feature recovery, support \\
    Ablation & real datasets except isolet & 22 & 8 & CIF ranking ablation \\
    Runtime & 9 real and 14 synthetic datasets & 15 & 8 & runtime ablations \\
    Benchmark-size & classification datasets with at least 500 observations except madelon; regression datasets with at least 200 observations & 9 & 5 & wide-dataset runtime check \\
    Mechanism-matched & cforest pairwise panel without the excluded RDC runs & 20 & 8 & CIF ranked by out-of-bag permutation importance \\
    Max-type extension & every real and synthetic benchmark dataset & 31 & 16 & max-type ranking runs \\
    Head-to-head & Benchmark-size classification datasets and madelon & 10 & 0 & forest timing of the two Stage~B rules \\
    Stage~B real & vowel, spam, page-blocks; imports-85, residential, facebook & 3 & 3 & threshold-test comparison \\
    \bottomrule
  \end{tabular}
  \caption{Dataset panels, with the number of datasets per task. The
  Benchmark-size panel is the set of performance datasets of
  \citet{MilletichDownesHirst2026citrees}. The max-type extension runs on every
  benchmark dataset and is ranked on the all-method panel.}
  \label{tab:panels}
\end{table}

\subsection{Datasets}
\label{sec:experiments-datasets}

The real-data benchmark includes 23 classification datasets and 8 regression
datasets. The separate high-$p$ analysis restricts attention to the
15 classification datasets and 6 regression datasets with more than 100
features. Feature recovery uses 8 synthetic classification datasets and 8
synthetic regression datasets with known informative features and controlled
noise structure. \Cref{app:dataset-inventory} lists every dataset and
summarizes the synthetic designs.

\subsection{Compared methods}
\label{sec:experiments-methods}

The benchmark includes three types of feature selectors: permutation-test
filters, embedded tree or ensemble methods, and wrapper baselines.
Permutation-test filters compute univariate dependence statistics and rank
features by their permutation $P$-values: multiple correlation (MC) and
randomized dependence coefficient (RDC) for classification, and Pearson
correlation (PC), distance correlation (DC), and RDC for regression. The RDC
keeps the empirical-CDF random feature map of \citet{LopezPaz2013RDC} but
replaces the canonical-correlation solve with the maximum absolute pairwise
correlation across 10 random projections. Distance correlation follows
\citet{Szekely2007DistanceCorrelation}. CIT and CIF use the same MC, PC, and
RDC statistics as Stage~A selectors.

Embedded methods rank features using tree or ensemble importance scores. These
methods are CART decision trees (DT), single randomized trees with random split
thresholds (RT), CIT, CIF, random
forests (RF) \citep{Breiman2001RandomForests}, ExtraTrees
\citep{Geurts2006ExtraTrees}, XGBoost \citep{ChenGuestrin2016XGBoost},
LightGBM \citep{Ke2017LightGBM}, CatBoost
\citep{Prokhorenkova2018CatBoost}, and the partykit reference implementations
ctree and cforest
\citep{RPartykitPackage,Hothorn2006UnbiasedRecursivePartitioning,HothornZeileis2015Partykit}.
Wrapper baselines score features through auxiliary random-forest fits: Boruta
\citep{Kursa2010Boruta}, permutation importance (PI) on a 20\% held-out split of
the training fold \citep{Breiman2001RandomForests,Strobl2007BiasRFVariableImportance},
conditional permutation importance (CPI)
\citep{DebeerStrobl2020ConditionalPermutationImportanceRevisited}, and
random-forest recursive feature elimination (RF-RFE)
\citep{Guyon2002RFE,Pedregosa2011ScikitLearn}. CPI permutes each feature within
five-bin percentile strata of the features whose absolute correlation with it
is at least 0.5, and globally when there are none. A further comparator ranks
the RF by out-of-bag permutation importance, which scores every training row.

Ranking rules use each method's natural score direction. Embedded tree and
ensemble baselines, PI, and CPI rank features by decreasing importance.
Permutation-test filters rank by increasing $P$-value, with dependence score as
the secondary key. The ctree reference ranks features by split-use counts.
cforest is ranked by partykit's out-of-bag permutation importance, the only
importance the package provides for forests, while CIF uses split importance, so
the two forests differ in importance mechanism as well as in how their trees
are grown. The mechanism-matched check (\cref{app:cif-perm-check}) ranks CIF by
cforest's own importance to separate the two. Boruta and RF-RFE rank by their
fitted rank variables, with smaller ranks first. Ties are broken by a
deterministic order.

\subsection{Configuration selection}
\label{sec:experiments-repro}

Several methods were evaluated under multiple configurations: four each for CIT
and CIF, five for XGBoost, four for cforest, and two each for LightGBM and
ctree. Permutation-test filters used one configuration per statistic. For each
task, the paper reports one configuration per method, selected by mean score
over real datasets, downstream models, and the standard values
$k \in \{5,10,25,50,100\}$. Methods with only one configuration are treated as
fixed. The selected configuration is then held fixed across datasets for that
task. The selected CIT configurations use the MC selector in classification and
the RDC selector in regression; the selected CIF configurations use RDC in both
tasks with 100 trees and square-root feature sampling, with honesty in
regression only (\cref{tab:selected-configs}).

Because configuration selection and reporting use the same real-data
benchmark, the reported scores, intervals, and rank summaries are descriptive.
As a sensitivity check, we repeat configuration selection in leave-one-dataset-out
(LODO) form for both tasks, on the complete-case and all-method panels: each
dataset is evaluated using configurations chosen from the other datasets of its
panel. The downstream models, $k$ values, and averaging within each dataset are
unchanged. For the synthetic analysis, each method is assigned the
configuration with the highest average recovery of the true informative
features over $k = 5, 10, 25, 50,$ and $100$.

\subsection{Ablations}
\label{sec:experiments-ablations}

The runtime ablation refits the selected CIT and CIF configurations on the
runtime panel with one control changed at a time: adaptive stopping, feature
scanning, threshold scanning, and feature muting are switched off; the two
corrections are removed together; the 256-bin histogram is replaced by exact
thresholds or, for CIF, by a 32-bin histogram; and, for CIF, bootstrap sampling
is switched off and the max-type rule replaces the per-threshold rule. Both
estimators are also refitted at the auto budget, with adaptive stopping and
with no stopping, because the minimum budget leaves the stopping rule no room
to act. Each dataset is fitted at five seeds, and every time is the second of
two identical fits, so compilation is excluded. On the Benchmark-size panel the
tree is refitted with adaptive stopping off, with threshold scanning off, and
without the two corrections, and the forest with adaptive stopping off.

To isolate which parts of the CIF ranking matter, a ranking ablation on the
ablation panel holds the selected CIF configuration fixed and changes one
component at a time: replacing fitted split importance with split counts,
disabling feature muting, disabling bootstrap sampling, or reducing the forest
to one tree. Each variant is evaluated with the same downstream learners,
folds, seeds, and standard top-$k$ values as the main benchmark, paired with the
reference configuration within seed and fold.

\subsection{Stage~B comparison}
\label{sec:experiments-stageB}

This study compares the per-threshold rule and the max-type rule under the MC
and PC selectors, with adaptive stopping and with no stopping, at five seeds
unless stated otherwise. Both rules run at the auto budget, so adaptive
stopping can act; the max-type rule uses 999 permutations. Candidate thresholds
come from the histogram method at 16, 64, and 256 bins, capped at the number of
distinct values in the node. The study has five parts.

The null calibration estimates the false-split rate of a depth-one tree with
five independent Gaussian predictors and an independent response, at
$n\in\{100,200,500\}$, with 2,000 replicates per cell and seed. The rate covers
the whole node decision, Stage~A followed by Stage~B.

The Stage-B-only study measures the threshold test by itself. It uses one
Gaussian predictor, sets the Stage~A level to 1, and fits a depth-one tree at
nominal Stage~B levels from 0.02 to 0.50, with 500 null replicates per cell,
level, and seed over three seeds. It repeats the planted step of the power part
with 200 replicates and compares power at matched realized size: within each
cell, power is interpolated to a common size of 0.05, or to the largest size
both rules reach. Two reruns check that comparison: a fourth seed with 1,500
null and 1,000 alternative replicates, and a design that moves the planted step
from the median to the 85th percentile of the predictor.

The scaling part times one tree on ten Gaussian predictors with a planted linear
signal in three of them, at $n\in\{1{,}000,4{,}000,16{,}000\}$, as the median of
three fits with a 1,200-second cap.

The real-data part fits one tree on the six Stage~B real datasets at 256 bins
with five-fold cross-validation, recording held-out score, tree depth, fitting
time, and the Spearman agreement of the two rules' feature importances.

The power part fits a depth-one tree on five Gaussian predictors whose response
depends on the first through a step at zero, with class-probability shift
$\delta\in\{0.05,0.10,0.20\}$ in classification and mean shift
$\delta\in\{0.2,0.4,0.8\}$ in regression, with 500 replicates per cell and seed.
It records the rate of any split, the rate of a split on the planted predictor,
and the error of the chosen threshold.

\subsection{Head-to-head timing and max-type extension}
\label{sec:experiments-h2h}
\label{sec:experiments-maxt}

The head-to-head timing fits 100-tree classification forests in the
configuration recommended by \citet{MilletichDownesHirst2026citrees}, once with
each Stage~B rule. That configuration uses the MC selector, the minimum budget,
the two corrections, scanning, muting, the 256-bin histogram, and depth at most
3. The cells are the ten head-to-head datasets and a synthetic grid of Gaussian
predictors with a weak or strong planted signal, from $n=500$ to $32{,}000$ and
$p=20$ to $500$.

The max-type extension reruns the four CIT and four CIF candidate
configurations with the max-type rule in place of the per-threshold rule and
nothing else changed, on every benchmark dataset at the five benchmark seeds.
Evaluation, configuration selection, and ranking are unchanged, with the
max-type family standing in for its per-threshold counterpart.

\subsection{Mechanism studies}
\label{sec:experiments-mechanism}

Three studies examine why the forest rankings behave as they do. The
feature-use study uses $n=250$, $p \in \{100, 500, 1{,}000\}$, and 1, 2, 5, or
10 informative features, and compares CIT, DT, RT, CIF, CIF-all, RF, and
ExtraTrees, with 1,000 trees for the forests. The weak-signal cardinality
designs place 10 informative Gaussian features among 25 noise features with 500
levels, 25 binary noise features, and 10 Gaussian noise features, at $n=1{,}000$
and signal strengths 0.05, 0.10, and 0.20; CIF and CIT use the MC selector in
classification and the PC selector in regression. The support measurement
counts, on each synthetic dataset over 25 fits, the features that win at least
one split, and how many of them are informative.
